\cotefile{63}
\scanpage{87}\mspage{T-1}

\begin{center}
COURBES ELLIPTIQUES - FORMULAIRE

(Traduction libre d'un manuscrit inédit de TATE)
\end{center}

\figure{\begin{tikzpicture}[x=0.016cm,y=0.016cm]
\draw[gray,thin] (0,0) .. controls (8,0) and (13,9) .. (22,14) .. controls (35,29) and (56,42) .. (77,43) .. controls (93,43) and (95,46) .. (97,50) -- (101,46);
\draw[gray,thin] (25,15) .. controls (32,0) and (43,-13) .. (47,-27);
\end{tikzpicture}}

1. \uline{Formules générales}

Si l'on part d'une c\deleted{\unsure{a|u}}ourbe elliptique $E$, munie d'une origine $\unsure{O|0}$, et d'une forme de première espèce $\pi\ne 0$, on peut trouver $x$ et $y$ sur $E$, telles que $x$ (resp. $y$) ait un pôle double (resp. triple) en $O$, et soit morphique ailleurs; de plus, si $\pi=dt+\ldots$, on peut s'arranger (quitte à multiplier $x$ et $y$ par des constantes) pour que $x\handwritten{=}{\unsure{t|+}}^{-2}+\ldots$, $y=-t^{-3}+\ldots$

Dans le plongement projectif défini par le diviseur $3(O)$, l'équation de la courbe $E$ s'écrit :
{\makeatletter\tagsleft@true\makeatother
\[
 y^{2}+a_{1}xy+a_{3}y=x^{3}+a_{2}x^{2}+a_{4}x+a_{6}
 \tag*{$(\ast)$}
\]
}

(Homogénéite : $y$ est de poids $3$, $x$ de poids $2$, et les $a_{i}$ de poids $i$. Cela signifie que, si l'on remplace $\pi$ par $u\pi$, $x$ est remplacé par $u^{-2}x$, $y$ par $u^{-3}y$, etc.)

Inversement, partons d'une équation de la forme $(\ast)$. On posera :
\[
\begin{aligned}
& b_{2}=a_{1}^{2}+4a_{2} \qquad b_{4}=a_{1}a_{3}+2a_{4} \qquad b_{6}=a_{3}^{2}+4a_{6} \\
& -b_{8}=a_{1}a_{3}a_{4}+a_{4}^{2}-a_{1}^{2}a_{6}-a_{2}a_{3}^{2}-4a_{2}a_{6} \\
& b_{8}=b_{2}a_{6}+a_{2}a_{3}^{2}-a_{1}a_{3}a_{4}-a_{4}^{2}=a_{2}b_{6}+a_{1}^{2}a_{6}-a_{4}^{2}-a_{1}a_{3}a_{4} \\
& \hspace{4em}-4\,b_{8}=b_{4}^{2}-b_{2}b_{6}
\end{aligned}
\]
% archive sheet number: 86

\scanpage{88}\mspage{T-2}

On pose également:
\[
 c_{4}=b_{2}^{2}-24\,b_{4}\qquad -c_{6}=b_{2}^{3}-36b_{2}b_{4}+216b_{6}
\]

(Les $c_{i}$ sont liés aux classiques $g_{2}$ et $g_{3}$ par les formule\rlap{\hspace{-.16em}\handwritten{\tikz[baseline=0pt]{\draw[line width=.32pt] (0,.04ex) .. controls (.18em,-.30ex) and (.45em,-.52ex) .. (.73em,-.71ex);}}}s:
\[
 c_{4}=\underset{\handwritten{2^{\unsure{2|2'}}\mathord{.}3}}{12}\,g_{2}
 \qquad c_{6}=\underset{\handwritten{2^{3}\mathord{.}3^{3}}}{216}\,g_{3}\ .)
\]

Et :
\[
\begin{aligned}
\handwritten{\Delta}&=-b_{2}^{2}b_{8}-8b_{4}^{3}-27b_{6}^{2}+9b_{2}b_{4}b_{6}\\
&=b_{4}^{3}-27b_{6}^{2}+b_{8}(36b_{4}-b_{2}^{2}).
\end{aligned}
\]

On a :
\[
 \underset{\handwritten{2^{6}\,3^{3}}}{1728}\,\handwritten{\Delta}=c_{4}^{3}-c_{6}^{2}\ .
\]

On pose enfin (pour $\handwritten{\Delta}$ inversible) :
\[
 j=\frac{c_{4}^{3}}{\handwritten{\Delta}}\qquad
 (j-1728)\handwritten{\Delta}=c_{6}^{2}\qquad
 \frac{j}{j-1728}=\frac{c_{4}^{3}}{c_{6}^{2}}
\]

Si l'on pose $F=y^{2}+a_{1}xy+a_{3}y-x^{3}-a_{2}x^{2}-a_{4}x-a_{6}$, on a :
\[
 F_{y}=2y+a_{1}x+a_{3}\qquad -F_{x}=3x^{2}+2a_{2}x+a_{4}-a_{1}y\ .
\]

Si $(\ast)$ définit bien une courbe elliptique (i.e. si $\handwritten{\Delta}$ est inversible), on peut écrire la forme différentielle $\pi$ de départ sous la forme $\pi=dx/F_{y}=-dy/F_{x}$, i.e. :
\[
 \pi=\frac{dx}{2y+a_{1}x+a_{3}}=\frac{dy}{3x^{2}+2a_{2}x+a_{4}-a_{1}y}
\]
% archive sheet number: 87

\scanpage{89}\mspage{T-3}

2. \uline{Changements de coordonnees}

Soit $C'=y'{}^{2}+a'_{1}x'y'+\ldots$ une autre courbe elliptique, et soit $f:C'\mathrel{\handwritten{\to}}C$ un isomorphisme transformant $\unsure{O|0}$ en $O$. On a alors :
\[
\begin{aligned}
x\circ f&=u^{2}x'+r\\
y\circ f&=u^{3}y'+su^{2}x'+t\ ,
\end{aligned}
\qquad\qquad \pi\unsure{\mathord{.}|\circ}f=u^{-1}\pi'
\]
où $u,r,s,t$ sont des constantes, $u$ étant inversible.

Les coefficients $a'_{i}$ sont lies aux $a_{i}$ par les formules :
\[
\begin{aligned}
ua'_{1}&=a_{1}+2s\\
u^{2}a'_{2}&=a_{2}-sa_{1}+3r-s^{2}\\
u^{3}a'_{3}&=a_{3}+ra_{1}+2t\\
u^{4}a'_{4}&=a_{4}-sa_{3}+2a_{2}-(t+rs)a_{1}+3r^{2}-2st\\
u^{6}a'_{6}&=a_{6}+ra_{4}+r^{2}a_{2}+r^{3}-ta_{3}-t^{2}-rta_{1}
\end{aligned}
\]

Pour les $b'_{i}$ :
\[
\begin{aligned}
u^{2}b'_{2}&=b_{2}+12r\\
u^{4}b'_{4}&=b_{4}+rb_{2}+6r^{2} &&\text{(\unsure{équation|equation} de discr. }c_{4}\text{)}\\
u^{6}b'_{6}&=b_{6}+2rb_{4}+r^{2}b_{2}+\text{\uline{$4$}}r^{3}
    &&\text{(équ.de discr. }16\,\handwritten{\Delta}\text{ )}\\
u^{8}b'_{8}&=b_{8}+3rb_{6}+3r^{2}b_{4}+r^{3}b_{2}+\text{\uline{$3$}}r^{4}
\end{aligned}
\]

Pour les $c'_{i}$ :
\[
\begin{aligned}
 &u^{4}c'_{4}=c_{4}&&u^{6}c'_{6}=c_{6}\\
\text{et}\quad &u^{12}\handwritten{\Delta}'=\handwritten{\Delta}&&j'=j
\end{aligned}
\]
% archive sheet number: 88

\scanpage{90}\mspage{T-4}

3. \uline{Forme canonique} ( $p\ne 2,3$)

C'est la forme
\[
 y^{2}=x^{3}+a_{4}x+a_{6}\qquad\qquad \pi=dx/2y
\]

On a :
\[
 c_{4}=-48\,a_{4}\ ,\quad c_{6}=-4\mathord{.}216\,a_{6}\ ,\quad
 \handwritten{\Delta}=16(-4a_{4}^{3}-27\,a_{6}^{2}).
\]

En posant $Y=2y$, $X=x$ , on retombe sur la forme de Weierstrass
\[
 Y^{2}=4X^{3}-g_{2}X-g_{3}\qquad\qquad \pi=dX/Y
\]
avec
\[
\begin{aligned}
g_{2}&=-4a_{4}=\frac{1}{12}\,c_{4} & g_{3}&=-4a_{6}=\frac{1}{216}\,c_{6}\\
\handwritten{\Delta}&=g_{2}^{3}-27g_{3}^{2} & \unsure{j|J}&=1728\,g_{2}^{3}/\handwritten{\Delta}\ .
\end{aligned}
\]

Les isomorphismes d'une telle courbe sur une autre sont donnés simplement par :
\[
 x\unsure{\circ|\mathord{.}}f=u^{2}x'\qquad\qquad y\;f=u^{3}y'\ .
\]

Deux courbes ayant le même $j$ sont isomorphes sur la clôture séparable de $k$ (corps de base - si c'est un schéma de base, il faut en prendre un revêtement). Plus précisément :

\uline{Cas général} : $j\ne 0,1728$\quad (i.e. $a_{4}\ne 0$, $a_{6}\ne 0$) .

Les deux courbes $C$ et $C'$ sont isomorphes si et seulement si $j=j'$ et $a_{4}a'_{6}/a'_{4}a_{6}$ est un carré ; le plus petit corps sur lequel $C$ et $C'$ deviennent isomorphes est le corps $k'$ obtenu en adjoignant à $k$ la racine carrée de $a_{4}a'_{6}/a'_{4}a_{6}$ .

Les automorphismes de $C$ correspondent à $u=\pm 1$.

\uline{Cas}\quad $j=1728$\quad (i.e. $a_{6}=0$) .

Ici $C\approx C'$ équivaut à $j=j'$ et $a_{4}/a'_{4}\in k^{\ast 4}$ .

Les automorphismes sont donnés par les $u$ tels que $u^{4}=1$.

Courbe type : $y^{2}=x^{3}-x$ .
% archive sheet number: 89

\scanpage{91}\mspage{T-5}

\uline{Cas}\quad $j=0$\quad (i.e. $a_{4}=0$) .

Ici $C\approx C'$ équivaut à $j=j'$ et $a_{6}/a'_{6}\in k^{\ast 6}$ .

Les automorphismes sont donnés par les $u$ tels que $u^{6}=1$.

Courbe type : $y^{2}=x^{3}-1$.

4. \uline{Forme canonique} ( $p=3$)

Du fait que $2$ est inversible, on peut tout d'abord écrire l'équation sous la forme :
\[
 y^{2}=x^{3}+a_{2}x^{2}+a_{4}x+a_{6}=f(x).
\]

On a :
\[
\begin{aligned}
 &b_{2}=a_{2}\ ,\quad b_{4}=-a_{4}\ ,\quad b_{8}=a_{4}^{2}-a_{2}a_{6}
 \qquad\quad \pi=-dx/y\\
 &c_{4}=a_{2}^{2}\ ,\quad c_{6}=a_{2}^{3}\ ,\quad
 \handwritten{\Delta}=a_{2}^{2}a_{4}^{2}-a_{2}^{3}a_{6}-a_{4}^{3}\ .
\end{aligned}
\]

Ici encore, $\handwritten{\Delta}$ est, à un facteur inversible près, le discriminant de $f$. Le fait que $C$ soit elliptique se traduit donc par $\handwritten{\Delta}\ne 0$, et $j=a_{2}^{6}/\handwritten{\Delta}$ . Deux courbes $C$ et $C'$ de même $j$ sont isomorphes après extension séparable de $k$. Plus précisément:

\uline{Cas général} : $j\ne 0$\quad (i.e. $a_{2}\ne 0$).

On peut faire disparaître le terme en $x$. Cela donne la forme réduite :
\[
 y^{2}=x^{3}+a_{2}x^{2}+a_{6}\ ,\quad
 \handwritten{\Delta}=-a_{2}^{3}a_{6}\ ,\quad j=-a_{2}^{3}/a_{6}\ .
\]

Un isomorphisme $f:C'\mathrel{\handwritten{\to}}C$ est donné par
\[
 x\;f=u^{2}x'\qquad y\unsure{\mathord{.}|\circ|\phantom{\circ}}f=u^{3}x'\ .
\]

Pour que $C\approx C'$ il faut et il suffit que $j=j'$ et que $a_{2}/a'_{2}\in k^{\ast 2}$ . Les automorphismes de $C$ correspondent à $u=\pm 1$.
% archive sheet number: 90

\scanpage{92}\mspage{T-6}

\uline{Cas}\quad $j=0$\quad (i.e. $a_{2}=0$).

Forme réduite :
\[
 y^{2}=x^{3}+a_{4}x+a_{6}\qquad \handwritten{\Delta}=-a_{4\unsure{.|\phantom{.}}}^{3}
 \qquad \pi=a_{4}^{-1}dy
\]

Isomorphismes :
\[
 x\unsure{\mathord{.}|\circ}f=u^{2}x'+r\ ,\quad
 y\mathord{.}f=u^{3}y'
\]

avec :
\[
 u^{4}a'_{4}=a_{4}\ ,\quad u^{6}a'_{6}=a_{6}+ra_{4}+r^{3}
\]

Deux courbes $C$ et $C'$ de même $j$ sont isomorphes si et seulement si :
\[
 a_{4}/a'_{4}\in k^{\ast 4}\qquad\text{et}\quad
 (a_{4}/a'_{4})^{3/2}a'_{6}-a_{6}\quad\text{est de la forme }r^{3}+ra_{4}.
\]

Cette condition est en tout cas vérifiée sur une extension séparable $k'/k$ de degré $\handwritten{\leqslant}\,12$.

Les automorphismes de $C$ sont donnés par :
\[
\begin{aligned}
u&=\pm 1 & r^{3}+a_{4}r&\deleted{\unsure{+(a_{6}-a'_{6})|+a_{6}-a'_{6}}}=0\\
u&=\pm i & r^{3}+a_{4}r&+\deleted{\unsure{(|\phantom{(}}}\handwritten{2}a_{6}\deleted{\unsure{+a'_{6})|-a'_{6})|+a'_{6}|-a'_{6}}}=0\ .
\end{aligned}
\]

Ils forment un groupe d'ordre $12$ (sur la clôture séparable de $k$).

L'anneau d'endomorphismes est un ordre maximal d'un corps de quaternions ramifié en $3$ et $\handwritten{\infty}$ .

Courbe type :
\[
 y^{2}=x^{3}-x.
\]

Dans ce cas, les automorphismes sont donnés par :
\[
 u^{4}=1\ ,\quad r^{3}-r=0\quad\text{(i.e. }r\in\text{\uuline{F}}_{3}\text{)}.
\]
% archive sheet number: 91

\scanpage{93}\mspage{T-7}

5. \uline{Forme canonique} ( $p=2$)

On a ici $ua'_{1}=a_{1}$ (cf. p.T-3). On distingue deux cas :

\uline{Cas général} : $a_{1}\ne 0$ (equivaut à $j\ne 0$).

En choisissant convenablement $u,r,t$, on s'arrange pour que $a_{1}=1$, $a_{3}=0$, $a_{4}=0$. D'où la forme canonique :
\[
 y^{2}+xy=x^{3}+a_{2}x^{2}+a_{6}\qquad\qquad \pi=dx/x\ .
\]

On a :
\[
 b_{2}=1,\ b_{4}=b_{6}=0,\ b_{8}=a_{6},\ c_{4}=1,\ \handwritten{\Delta}=a_{6}\ ,\
 j=\frac{1}{a_{6}}\ .
\]

Isomorphismes :
\[
 x\circ f=x'\ ,\quad y\circ f=y'+sx'\ ,
\]
avec :
\[
\begin{aligned}
 a'_{2}&=a_{2}+s^{2}-s\\
 a'_{6}&=a_{6}\ .
\end{aligned}
\]

Deux courbes $C$ et $C'$ de même $j$ sont isomorphes si et seulement si $a'_{2}-a_{2}$ est de la forme $s^{2}-s$ ; cette condition est vérifiée sur une extension séparable de $k$ de degré $\handwritten{\leqslant}2$.

Le groupe d'automorphismes de $C$ a deux éléments, correspondant à $s=0,1$.

Courbe type :
\[
 y^{2}+xy=x^{3}+a_{6}\ .
\]

\uline{Cas}\quad $a_{1}=0$\quad (équivaut à $j=0$).

En choisissant convenablement $r$, on s'arrange pour que $a_{2}=0$. D'où la forme canonique :
\[
 y^{2}+a_{3}y=x^{3}+a_{4}x+a_{6}\qquad\qquad \pi=dx/a_{3}\ .
\]

On a :
\[
 b_{2}=b_{4}=0,\ b_{6}=a_{3}^{2},\ b_{8}=a_{4}^{2}\ ,\
 \handwritten{\Delta}=a_{3}^{4}\ ,\ j=0.
\]
% archive sheet number: 92

\scanpage{94}\mspage{T-8}

Isomorphismes :
\[
\begin{aligned}
 u^{3}a'_{3}&=a_{3}\\
 u^{4}a'_{4}&=a_{4}+sa_{3}+s^{4}\\
 u^{6}a'_{6}&=a_{6}+s^{2}a_{4}+ta_{3}+s^{6}+t^{2}\ .
\end{aligned}
\]

Deux courbes $C$ et $C'$ de même $j$ ( égal à $0$) sont isomorphes si les \unsure{équations|equations} ci-dessus sont resolubles en $u,s,t$, ce qui est le cas sur une extension séparable de $k$ de d\textsuperscript{o} $\handwritten{\leqslant}\,24$.

Courbe type :
\[
 y^{2}-y=x^{3}\ .
\]

Son groupe d'automorphismes (sur la clôture séparable de $k$) est d'ordre $24$ ; il est donné par :
\[
 u^{3}=1\ ,\quad s^{4}+1=0,\quad t^{2}+t+s^{3}+s^{2}=0.
\]

L'anneau d'endomorphismes de $C$ est isomorphe à l'anneau des ``quaternions entiers'' de Hurwitz : ensemble des $a+bi+cj+dk$ où les $a,b,c,d$ sont simultanément, soit entiers, soit demi-entiers (i.e. entier $+1/2$).
% archive sheet number: 93

